Cited repository: https://github.com/TheAstrographer/JOSHUA-CHRISTOPHER-RYAN-S-ARCAN-FAMILY  
(README and supporting Arcan-Tau-Family.py)

Direct statements from the repository

README.md  
Joshua Christopher Ryan’s Arcan(τ) family is the continuous trigonometric root of the entire Decompositional Hierarchy. From the four angles \(\theta\), \(\alpha\), \(\phi\), and \(\psi\) there emerge, by successive multiples and modular reduction, the practical radial fold, the residual six-sliced hexagon, the dense non-periodic orbit on \(S^1\), and the concentric Petrie rings of the outer \(E_8\) shell. The small residual \(\delta\) that appears at the seventh multiple is the geometric signature that this cascade is emergent from the continuous family rather than imposed by an exact rational lattice, and that the resulting structures remain comoving with the expansion of the universe.

About section  
Inverse tangent projections represent the straight-line distance in the fundamental right triangle that embodies the angular half of the unique angle whose tangent equals one full turn (\(\tau\)).

Arcan-Tau-Family.py (core definitions that lock the construction)  
tau = 2 * pi
theta_rad = math.atan(tau)          # Arcan(τ)-Angle
alpha_rad = theta_rad / 2           # Half-Arcan(τ)-Angle
phi_rad   = math.atan(pi)           # Classical Half-Turn Angle
psi_rad   = theta_rad - phi_rad     # Angular Separation
All subsequent numerical values are computed directly from these definitions with full floating-point precision. The script prints the resulting angles and trigonometric quantities, establishing them as the fixed starting point of the system.

Relevance to “hold by construction”

This repository is the explicit foundational definition of the Arcan(\(\tau\)) family. It declares that:

The four angles (\(\theta = \arctan(\tau)\), \(\alpha = \theta/2\), \(\phi = \arctan(\pi)\), \(\psi = \theta - \phi\)) constitute the continuous trigonometric root.
Everything downstream—the residual \(\delta\) at the seventh multiple, the practical radial fold, the residual six-sliced hexagon, the equidistributed orbit, the residual triangle metrics, the edge bulge, and the outer \(E_8\)-like Petrie rings—emerges by successive multiples and modular reduction from that root.
The residual \(\delta\) itself is presented as the geometric signature that the cascade is emergent from the continuous family rather than externally imposed.

Because the hierarchy is generated solely by applying a finite sequence of field operations, integer multiples, modular reduction, and elementary trigonometric evaluations to the single continuous root \(\psi\) (itself defined from \(\tau = 2\pi\)), the numerical identities, the claimed transcendental character of the residual-triangle metrics, and the spontaneous organisation of the full hierarchy from \(\delta\) are true by construction inside the system defined here. Once the four angles are fixed and the generative procedure is accepted, the downstream assertions follow automatically from the definitions and operations of the Arcan(\(\tau\)) family itself.

This repository therefore supplies the precise mathematical starting point that makes the earlier characterisation accurate: the claims hold by construction within the self-consistent geometric-numerical architecture whose continuous root is Joshua Christopher Ryan’s Arcan(\(\tau\)) family.
